\chapter{案例分析与应用讨论}

\section{案例一：火灾疏散引导}
假设停车场 B2 层发生火情，需要实时定位疏散人员。本文方法可持续输出位置和可信度，指挥端可据此判断人员是否偏离疏散通道。

\section{案例二：救援队员协同定位}
在多名救援队员同时作业场景下，系统可对每个终端独立定位。若某队员轨迹出现异常跳变，可结合可信度告警快速核查设备状态。

\section{案例三：局部信标失效}
模拟部分区域停电导致 30\% 信标失效。对比结果显示：仅靠 RSSI 反演会出现明显漂移，而本文方法依托运动约束仍能维持连续轨迹。

\section{调度价值分析}
定位结果在应急系统中可支持：
\begin{enumerate}
  \item 疏散路径动态推荐；
  \item 被困人员最后位置回放；
  \item 救援任务区域覆盖评估；
  \item 风险区域二次进入预警。
\end{enumerate}

\section{落地限制}
本方案仍受以下条件限制：
\begin{itemize}
  \item 终端需持续扫描蓝牙，存在能耗开销；
  \item 高密度车辆金属反射会放大 RSSI 波动；
  \item 地图质量不足会削弱约束有效性。
\end{itemize}

\section{实践建议}
建议在部署前完成三项准备：信标健康巡检、路径损耗现场标定、车道拓扑电子化建模。三项准备充分时，定位结果稳定性可显著提升。
